% Generated draft: requires mathematical and source review.
\documentclass[11pt]{article}
\usepackage[margin=25mm]{geometry}
\usepackage{fontspec}
\setmainfont{Latin Modern Roman}
\usepackage{amsmath,amssymb,mathtools,unicode-math}
\setmathfont{Latin Modern Math}
\usepackage{enumitem,needspace,xurl,hyperref}
\setlength{\parindent}{0pt}
\setlength{\parskip}{5pt}
\setlength{\emergencystretch}{3em}
\newcommand{\N}{\mathbb{N}}
\newcommand{\Z}{\mathbb{Z}}
\newcommand{\Q}{\mathbb{Q}}
\newcommand{\R}{\mathbb{R}}
\newcommand{\C}{\mathbb{C}}
\newcommand{\F}{\mathbb{F}}
\newcommand{\A}{\mathbb{A}}
\newcommand{\PP}{\mathbb{P}}
\newcommand{\E}{\mathbb{E}}
\newcommand{\eps}{\varepsilon}
\newcommand{\dd}{\,\mathrm{d}}
\begin{document}
\section*{UnsolvedMath ID: 30005561}
\textbf{Source problem code:} catalogue entry
\subsection{Definitions and mathematical statement}
Let \(X\) be a smooth projective hyperk\"ahler variety, and let \(f\colon X\to B\) be a Lagrangian fibration: its general smooth fiber is Lagrangian for the holomorphic symplectic form on \(X\). Write \(d=\dim_{\C}X\); a Lagrangian fiber has dimension \(d/2\). Let \(L\) denote the first Chern class of the line bundle said in the supplied record to induce this fibration. The original passage refers to \(X,L\) ``as above,'' but the preceding hypotheses specifying this relationship are not supplied.

Write \(\operatorname{CH}^{r}(X)_{\Q}\) for the Chow group of codimension-\(r\) algebraic cycles modulo rational equivalence, with rational coefficients, and let \(c_{2j}(X)=c_{2j}(T_X)\in\operatorname{CH}^{2j}(X)_{\Q}\). A weighted homogeneous polynomial of weighted degree \(2i\) in the even Chern classes means a finite expression
\[
 c_I(X)=\sum_{\alpha}a_{\alpha}\prod_{j\geq1}c_{2j}(X)^{\alpha_j},
 \qquad a_{\alpha}\in\Q,\qquad
 \sum_{j\geq1}2j\alpha_j=2i,
\]
where each \(\alpha_j\) is a nonnegative integer. Thus \(c_I(X)\in\operatorname{CH}^{2i}(X)_{\Q}\). Products and powers below are taken in the Chow ring.

The supplied Conjecture 2 asks that, for every such polynomial and every intended admissible pair of integer indices \(i,k\) satisfying \(i\geq k\),
\[
 c_I(X)L^{n+1-k}=0
 \quad\text{in }\operatorname{CH}^{2i+n+1-k}(X)_{\Q}.
\]
Here \(i\geq0\) and \(n+1-k\geq0\) are necessary for the displayed polynomial degree and ordinary power to be defined. The complete intended index range is not supplied.

\emph{Unresolved dimension convention.} The clean statement and the opening of the quoted passage explicitly call \(n\) the dimension of \(X\). However, the same passage describes the general fiber class as proportional to \(L^n\), and assigns \(n=5\) to the LSV example. These contextual uses suggest that \(n\) instead denotes half the dimension of \(X\). Consequently, the intended hypothesis may be \(d=2n\), rather than the explicitly written \(d=n\). This draft retains the exponent \(n+1-k\) exactly and does not select a corrected dimension convention. A definitive formulation requires resolving this discrepancy and recovering the omitted hypotheses on \(L\) and the index range.

The source reports the following partial results, subject to the same unresolved notation and hypotheses. Its Theorem 2 asserts
\[
 c_{2i}(X)L^{n+1-k}=0\qquad(i\geq k).
\]
It also reports, for an LSV manifold obtained as a compactification of the intermediate Jacobian fibration associated to a cubic fourfold, equipped with its canonical Jacobian fibration,
\[
 c_2(X)^2L^{n-1}=0,\qquad n=5.
\]
Here the extracted typography \(c22\) is read as \(c_2(X)^2\), consistently with the passage's description of an additional Chern monomial. The source explicitly distinguishes the individual Chern-class relations from the conjecture for all higher monomials.

For context, the passage describes the individual Chern-class result as formal, proved by induction on \(i\), with the \(k=1\) case using the general abelian fiber and its trivial tangent and normal bundles. It describes the LSV monomial result as requiring additional geometry: the discriminant hypersurface is the dual hypersurface of the cubic fourfold \(Y\subset\PP^5\), inside \((\PP^5)^*\), and is reported to be rationally connected with \(\operatorname{CH}_0=\Z\). The passage explicitly says this property does not hold in general.

The source further reports rank loci \(X_{\leq k}\) and \(X_k\), where the fibration has rank at most \(k\) and exactly \(k\), respectively. With \(B_k=f(X_k)\), its Theorem 3 states
\[
 \dim B_k\leq k,\qquad \dim X_k/B_k\geq k.
\]
The precise rank convention, hypotheses on the base, and meaning of the latter relative-dimension notation are not specified in the supplied excerpt. These are recorded contextual claims, not additional conclusions established by this draft.
\subsection{Short English statement}
The conjecture asks for rational Chow-ring vanishing of every weighted homogeneous polynomial in the even Chern classes after multiplication by a specified power of the fibration class. The supplied record reports partial results for individual even Chern classes and for \(c_2(X)^2\) in the LSV case. The meaning of \(n\), the full index range, and some preceding hypotheses remain unresolved, so a definitive corrected statement requires mathematical review.
\subsection{Sources}
\begin{itemize}
\item Supplied local UnsolvedMath registry record, permanent ID \(30005561\), exact source problem code \texttt{catalogue entry}, title ``Chow Vanishing for Lagrangian-Fibered Hyperk\"ahler Varieties.'' The record places the problem in the Oberwolfach Reports Open Problems collection, with proposed year 2023 and no named proposer.
\item \emph{Recent Trends in Algebraic Geometry} (2023), Oberwolfach Reports, report 27/2023, pp.\ 1485--1542, DOI 10.4171/owr/2023/27. Supplied report and source URL: \url{https://doi.org/10.4171/owr/2023/27}. The supplied excerpt labels the target Conjecture 2 and the contextual results Theorems 2 and 3.
\item Voisin, \emph{On Chern classes of Lagrangian fibered hyper-K\"ahler manifolds}. Supplied literature link: \url{https://arxiv.org/abs/2305.09396}. The supplied assessment attributes the individual even Chern-class vanishings and an additional LSV monomial vanishing to Voisin. The excerpt refers to results in its reference [10] and to the LSV construction in its reference [4]; their full numbered bibliography entries are not supplied.
\item Supplied literature assessment, checked exactly \texttt{2026-08-22}: status ``partially solved.'' It reports that the individual even Chern classes and one additional LSV monomial are covered, while remaining monomial cases are unsettled. Its listed literature sources are \url{https://arxiv.org/abs/2305.09396} and \url{https://doi.org/10.4171/owr/2023/27}. This is supplied dated curation, not an independently verified literature assessment.
\item The registry labels the statement \texttt{corrected\_verified} and says it was reconstructed and checked against the official report. That verification is a supplied registry assertion; this draft does not independently verify it or claim to have visited either source.
\end{itemize}
\end{document}
